% Generated draft: requires mathematical and source review.
\documentclass[11pt]{article}
\usepackage[margin=25mm]{geometry}
\usepackage{fontspec}
\setmainfont{Latin Modern Roman}
\usepackage{amsmath,amssymb,mathtools,unicode-math}
\setmathfont{Latin Modern Math}
\usepackage{enumitem,needspace,xurl,hyperref}
\setlength{\parindent}{0pt}
\setlength{\parskip}{5pt}
\setlength{\emergencystretch}{3em}
\newcommand{\N}{\mathbb{N}}
\newcommand{\Z}{\mathbb{Z}}
\newcommand{\Q}{\mathbb{Q}}
\newcommand{\R}{\mathbb{R}}
\newcommand{\C}{\mathbb{C}}
\newcommand{\F}{\mathbb{F}}
\newcommand{\A}{\mathbb{A}}
\newcommand{\PP}{\mathbb{P}}
\newcommand{\E}{\mathbb{E}}
\newcommand{\eps}{\varepsilon}
\newcommand{\dd}{\,\mathrm{d}}
\begin{document}
\section*{UnsolvedMath ID: 30003709}
\textbf{Source problem code:} catalogue entry
\subsection{Definitions and mathematical statement}
The supplied statement does not specify the coefficient ring of cohomology. Its accompanying literature assessment discusses rational cohomology, suggesting the rational formulation below; this interpretation requires confirmation.

For the standard complex-torus interpretation, let \(T_n=(\C^\times)^n\), where \(n\geq 1\). A central toric arrangement \(A\) is a finite collection of connected algebraic subtori of complex codimension one in \(T_n\). Equivalently, each member has the form
\[
H_a=\{(z_1,\ldots,z_n)\in T_n:z_1^{a_1}\cdots z_n^{a_n}=1\},
\]
where \(a=(a_1,\ldots,a_n)\in\Z^n\) is nonzero and primitive. Here primitive means \(\gcd(|a_1|,\ldots,|a_n|)=1\). Every member contains the identity of \(T_n\). Define the complement by
\[
M(A)=T_n\setminus\bigcup_{H\in A}H.
\]
A layer is a connected component of an intersection \(\bigcap_{H\in S}H\), for a subset \(S\subseteq A\); the empty intersection is \(T_n\). The layer poset \(L(A)\) consists of these layers, ordered by reverse inclusion:
\[
X\leq Y\quad\Longleftrightarrow\quad X\supseteq Y.
\]
For a specified commutative coefficient ring \(k\), write \(H^\bullet(M(A);k)\) for singular cohomology with its grading, unit, and cup product.

The reconstruction question, conditional on these conventions and the choice of \(k\), is whether, for every pair of central toric arrangements \(A\) and \(B\), possibly in different-dimensional complex tori,
\[
H^\bullet(M(A);k)\cong H^\bullet(M(B);k)
\quad\text{as graded unital }k\text{-algebras}
\quad\Longrightarrow\quad
L(A)\cong L(B)\quad\text{as posets}.
\]
The supplied literature assessment suggests taking \(k=\Q\), but the original question leaves coefficients unstated. The displayed implication is therefore a proposed precise interpretation, not a verified resolution of that omission. The source also does not explicitly specify the ambient category, whether defining characters must be primitive, or whether algebra isomorphisms must preserve grading. The connected complex-hypertorus convention above must be checked against the intended source definition; allowing disconnected character kernels could change the class of arrangements.
\subsection{Short English statement}
Can the cohomology algebra of the complement of a central toric arrangement recover, up to poset isomorphism, the connected components of all intersections and their inclusion relations? The supplied curation suggests rational coefficients, but the original statement does not specify coefficients or all arrangement conventions. The problem is labeled open in the supplied record; that status is not independently verified.
\subsection{Sources}
\begin{itemize}
\item Supplied local registry record: \texttt{data/unsolvedmath/problems.json}, permanent UnsolvedMath ID \(30003709\), source problem code \texttt{catalogue entry}, title ``Recovering Toric Arrangement Posets from Complement Cohomology.'' Collection: Oberwolfach Reports Open Problems. The record gives proposed year \(2018\) and does not name a proposer.
\item Supplied report citation: \emph{Topology of Arrangements and Representation Stability} (2019), Oberwolfach Reports, DOI 10.4171/owr/2018/2, pp.\ 43--123. \url{https://doi.org/10.4171/owr/2018/2}. The background identifies Oberwolfach Reports as reports of mathematical workshops and their open-problem discussions. The supplied citation's 2019 date and its 2018 report identification are retained without reconciliation.
\item Supplied literature source: Pagaria, \emph{Two Examples of Toric Arrangements}. \url{https://arxiv.org/abs/1804.05767}.
\item Supplied literature source: \emph{Oberwolfach Report 2/2018}. \url{https://doi.org/10.4171/owr/2018/2}.
\item Literature assessment dated exactly \texttt{2026-08-21}: the supplied curation labels the problem open and states that the layer poset determines the rational cohomology algebra, that integral counterexamples show limits of the converse framework, and that no result establishes recovery of the toric layer poset from the rational algebra. These are supplied dated curation claims, not independently verified mathematical or literature conclusions. The background explicitly requests independent verification of time-sensitive status and cited claims.
\item The registry labels the statement \texttt{corrected\_verified} and says its mathematical content was preserved while notation, typography, or omitted definitions were repaired for self-contained publication. This is supplied registry metadata, not independent verification of this draft or its conventions.
\end{itemize}
\end{document}
